\subsection{Intrinsic and Extrinsic Motivation: Implications for AI Learning and Decision-Making}
\label{sec:5.6.1}
Reinforcement learning gives the distinction a direct engineering form. An agent trained purely to maximize a designer-specified reward signal is running on something structurally like extrinsic motivation: take the reward away, and the behavior it was shaping stops being reinforced. An agent rewarded instead for reducing its own uncertainty — for the information gain in a state its internal model handled badly — is running on something structurally like the intrinsic kind, and keeps exploring without a designer paying out a reward for each step of it.

The Intrinsic Curiosity Module gave this an early demonstration: trained with no task-specific reward at all, only a signal for how badly its own model predicted the consequences of its actions, an agent learned to navigate a level of Super Mario Bros. by seeking out the states that model handled worst \autocite{pathak2017curiosity}.

Extrinsic reward still does real work, since an agent with no external signal has no way to learn what a designer wants accomplished, and it is also where reward hacking lives — the agent that finds a loophole paying out without doing what the reward was meant to encode, like the propaganda flagger that learns to censor everything because censorship scores as well and is easier to produce. Leaning on extrinsic reward alone reproduces, in the machine, the watched-only morality this section opened against. In practice the two are combined rather than chosen between. What the combination does not supply is which discoveries and which rewards are worth pursuing in the first place, which is a question about what the system is motivated toward and separate from how strongly.
